\begin{table}[H]
\centering
\label{tab:pi0_estimate_FF}
\begin{threeparttable}
\caption{Aggregate Skill Estimate: Proportion of True Zero-Alpha Active Funds -- FF Subperiod}
\begin{tabular}{lrrrl}
\toprule
Metric & Estimate & $N$ & $\lambda$ & Interpretation \\
\midrule
$\hat{\pi}_0$: Proportion of True Zero-Alpha Active Funds & 76.0\% & 805 & 0.5 & Heterogeneous skill; majority zero-alpha \\
\bottomrule
\end{tabular}
\begin{tablenotes}
\item The proportion of true zero-alpha funds ($\hat{\pi}_0$) is estimated following \textcite{Storey2002} and \textcite{BarrasScailletWermers2010}, using p-values from Newey--West $t$-tests on full-period \textcite{Carhart1997} four-factor alphas. The estimator is $\hat{\pi}_0 = |\{p_i > \lambda\}| \,/\, [N(1-\lambda)]$, where $\lambda = 0.5$ is the standard tuning parameter following \textcite{Storey2002}. Funds with p-values exceeding $\lambda$ are unlikely to have non-zero true alpha; the density of p-values in $(\lambda, 1]$ provides a conservative estimate of the zero-alpha proportion, bounded above at 1. Sample: actively managed funds ($N = 805$), Jan 1995--Sep 2006 (Fama--French 2010 subperiod), performance-comparison subsample per flagged\_funds.xlsx. Passive and Unknown-classified funds are excluded. Compare with \textcite{FamaFrench2010}, Table~III, which reports an analogous bootstrap-based assessment over Jan 1984--Sep 2006.
\end{tablenotes}
\end{threeparttable}
\end{table}
